\chapter{Introdução}

A transformação digital no setor de saúde tem avançado de forma acelerada nas últimas
décadas, impulsionada pela necessidade de otimizar recursos, ampliar o acesso ao atendimento
e melhorar a precisão diagnóstica. O conceito de Saúde 4.0 integra tecnologias como Internet
das Coisas, Big Data e Inteligência Artificial para criar sistemas de saúde mais preditivos e
personalizados. Nesse contexto, a interseção entre a Ciência da Computação e a Medicina tem
proporcionado avanços significativos na forma como diagnósticos e triagens são realizados.

No entanto, apesar da modernização dos equipamentos de exames e dos procedimentos
cirúrgicos, a etapa fundamental da prática médica, a anamnese, permanece majoritariamente
manual, analógica ou pouco estruturada. A entrevista inicial entre médico e paciente, momento
crucial para a coleta da história clínica, consome tempo valioso que poderia ser direcionado à
análise clínica aprofundada e à humanização do cuidado.

Este trabalho propõe uma solução de engenharia de software para otimizar esse gargalo,
utilizando Modelos de Linguagem de Grande Escala para interação em linguagem natural e
Reconhecimento de Entidades Nomeadas para estruturação técnica dos dados coletados
\cite{singhal2023,sang2003}.

\section{Contextualização}

O setor de saúde enfrenta um desafio global de sobrecarga e sustentabilidade. Segundo a
Organização Mundial da Saúde, projeta-se um déficit global de profissionais de saúde até 2030.
Esse cenário impõe pressão sobre os médicos em atividade, levando a burnout e aumento da
probabilidade de erros associados à fadiga.

Profissionais médicos dedicam parcela significativa de sua carga horária ao preenchimento de
prontuários. Estudos indicam que a documentação clínica pode consumir de 35\% a 40\% do
tempo total de uma consulta. Além disso, as queixas dos pacientes chegam aos sistemas de
saúde de forma difusa e desorganizada, exigindo esforço cognitivo intenso para filtrar sinais
relevantes.

Paralelamente a esse cenário, observa-se uma explosão na disponibilidade de dados e na
capacidade computacional. A evolução dos modelos generativos baseados em Transformer,
como as famílias GPT e Llama, permitiu que sistemas computacionais interpretassem e
gerassem linguagem natural com maior fluidez e precisão \cite{vaswani2017,metaai2024}.

Contudo, a aplicação direta dessas ferramentas na saúde esbarra em desafios críticos de
engenharia e ética: confiabilidade das informações geradas, privacidade de dados sensíveis e
interoperabilidade com sistemas legados.

\section{Definição do Problema}

O problema central abordado neste trabalho é a ineficiência na coleta e estruturação dos dados
de queixa principal dos pacientes na etapa de pré-consulta.

Atualmente, o fluxo de informação apresenta uma ruptura semântica entre o relato em texto
livre e o registro estruturado necessário para sistemas clínicos. A dependência de intervenção
humana para traduzir narrativas em dados padronizados gera:

\begin{itemize}
\item Gargalo operacional, reduzindo a capacidade de atendimento.
\item Perda de informação, com omissão de detalhes clínicos sutis.
\item Dados não padronizados, dificultando interoperabilidade e análise posterior.
\end{itemize}

A questão de pesquisa que norteia este trabalho é: como uma arquitetura de software híbrida,
combinando modelos generativos e extrativos, pode automatizar a anamnese prévia com
fidelidade dos dados clínicos extraídos?

\section{Objetivos}

\subsection{Objetivo Geral}

Desenvolver e validar uma arquitetura de aplicação web para realização de anamnese assistida,
capaz de interagir com o paciente em linguagem natural e entregar ao profissional de saúde um
relatório estruturado contendo sintomas, com abordagem baseada em microsserviços de IA.

\subsection{Objetivos Específicos}

\begin{itemize}
\item Realizar levantamento bibliográfico sobre arquiteturas de LLMs, técnicas de NER e
padrões de interoperabilidade na saúde.
\item Projetar e implementar uma API Backend for Frontend (BFF) em Node.js para
orquestrar inferências de forma segura e escalável.
\item Integrar e parametrizar o modelo Llama-3.1-8B-Instruct para condução do diálogo
clínico com técnicas de engenharia de prompt.
\item Integrar o modelo Biomedical-NER para extração automática de entidades.
\item Desenvolver uma interface reativa para validação do MVP.
\item Avaliar viabilidade técnica da solução por métricas qualitativas e capacidade de extração.
\end{itemize}

\section{Justificativa}

A relevância deste projeto se sustenta em três frentes.

No âmbito social, ferramentas de anamnese assistida podem ampliar eficiência em triagem e
telemedicina, especialmente em cenários de escassez de especialistas.

No âmbito técnico, a arquitetura híbrida proposta aumenta confiabilidade ao separar conversa
generativa da extração estruturada de dados, reduzindo impactos de alucinação factual.

No âmbito acadêmico, a solução documenta integração de ferramentas open-source em
arquitetura auditável e de baixo custo, contribuindo para pesquisas futuras em IA na saúde.

\section{Organização do Trabalho}

Este documento está organizado em cinco capítulos e apêndices. O Capítulo 1 apresenta
contexto, problema e objetivos. O Capítulo 2 traz a fundamentação teórica. O Capítulo 3
discute trabalhos relacionados. O Capítulo 4 detalha metodologia, arquitetura e resultados do
MVP. O Capítulo 5 apresenta considerações finais do TCC 1 e planejamento do TCC 2.
